\documentclass[12pt]{article}
\usepackage[utf8]{inputenc}
\usepackage[T1]{fontenc}
\usepackage{hyperref}
\title{Potigu\'ara Language Resources for K-12 Curriculum Development}
\author{Compiled by Claude Sonnet 5.5 (claude-sonnet-5-5) for Andr\'as Kornai}
\date{}
\begin{document}
\maketitle

\section{Introduction}
Potigu\'ara (ISO 639-3 \texttt{pog}) is a dormant Tupi-Guaran\'i language of the Brazilian northeast coast (Para\'iba, Rio Grande do Norte). The Potiguara ethnic population was 37,292 in the 2022 census, but the language is no longer a community first language. The 2022 census nevertheless recorded 806 people reported as speakers of Tupi-Potiguara, and sources mention revival efforts. Search found only secondary and ethnographic material, nothing like a curriculum.

\section{Primary Educational Resources}
No Potiguara primers or textbooks were found in search results.

\section{Linguistic and Reference Materials}
\begin{itemize}
\item Academic paper on Potiguara lexicon and language, UFPB journal (the search summary mentions mangrove vocabulary and language revival). \url{https://periodicos.ufpb.br/index.php/rics/article/download/74047/41582/230563}
\item Article in \emph{Tempus} (UnB) found in the Potiguara language search; its content was not verified in detail. \url{https://tempus.unb.br/index.php/tempus/article/download/1100/992/2369}
\end{itemize}

\section{Religious and Community Resources}
\begin{itemize}
\item A Uni\~ao (Para\'iba state newspaper), ``Para\'iba registrou 95 etnias ind\'igenas'': demographic context. \url{https://auniao.pb.gov.br/noticias/caderno_diversidade/paraiba-registrou-95-etnias-indigenas}
\end{itemize}

\section{Digital and Audio Resources}
None found.

\section{Curriculum Development Context}
Any curriculum would be a revitalisation or reclamation project, drawing on historical Tupi-Guaran\'i sources and the communities' own efforts. The Potiguara live in several indigenous lands, and Portuguese would remain the language of instruction. The sources above give background only.

\end{document}